% ===== 原文 pp.67--71 =====
% [原文 p.67]
\subsection{'t Hooft 反常与 SPT 相}

\subsubsection{'t Hooft 反常基础}

\para{纯 't Hooft 反常。} 当下面两点都成立时，$p$-形式对称性就存在纯 't Hooft 反常：
\begin{itemize}[leftmargin=1.4em,itemsep=1pt,topsep=1pt]
\item 背景场 $B_{p+1}$ 的规范变换不使配分函数保持不变
\begin{equation}
  Z[B_{p+1}]\neq Z[B^g_{p+1}]\,, \tag{4.34}
\end{equation}
其中 $B^g_{p+1}$ 与 $B_{p+1}$ 规范等价，即
\begin{equation}
\begin{gathered}
  B^g_{p+1}=B_{p+1}+\dd\Lambda_p\,,\qquad G^{(p)}\ \text{连续}\,,\\
  B^g_{p+1}=B_{p+1}+\delta\Lambda_p\,,\qquad G^{(p)}\ \text{离散}\,.
\end{gathered}
\tag{4.35}
\end{equation}
\item 即使向作用量中加入一个抵消项（它是 $B_p$ 的函数），仍然不可能使 $Z[B_{p+1}]=Z[B^g_{p+1}]$。
\end{itemize}

% [Tutor Note: p.67 第二个项目符号原文写 “a counter-term, which is a function of $B_p$”。抵消项应与所考察的 $p$-形式对称性背景场 $B_{p+1}$ 同变（见 (4.34) 中 $Z[B_{p+1}]$ 的变分），故此处 $B_p$ 应为 $B_{p+1}$；已用 700 dpi 裁图复核，原文确实印作 $B_p$，译文照录。]

\para{混合 't Hooft 反常。} 我们也会遇到这样的情形：有
\begin{equation}
  Z[B_{p+1}]=Z[B^g_{p+1}]\,, \tag{4.36}
\end{equation}
但如果我们同时打开另一个 $q$-形式对称性的背景场，则配分函数在规范变换下不再不变
\begin{equation}
  Z[B_{p+1},B_{q+1}]\neq Z[B^g_{p+1},B_{q+1}]\,, \tag{4.37}
\end{equation}
并且这种变化无法通过加入抵消项（现在它是 $B_{p+1}$ 与 $B_{q+1}$ 两者的函数）来修正。在这种情形下，我们说 $p$-形式与 $q$-形式对称性之间存在混合 't Hooft 反常。

也可能存在涉及两个以上对称性的混合 't Hooft 反常，其定义是上述情形的直接推广。

\begin{notebox}{其他类型的反常}
't Hooft 反常与另外两类反常密切相关，但不应当与它们混淆：
\begin{itemize}[leftmargin=1.4em,itemsep=1pt,topsep=1pt]
\item \para{ABJ 反常：} 当存在一个变换，它是经典拉格朗日量的对称性，却不是完整量子配分函数（其中还包括路径积分的测度）的对称性时，就会出现这类反常。这样的变换\textbf{不是} QFT 的整体对称性。人们常说，这样的变换是经典理论的对称性，但受到 ABJ 反常的困扰，因此它不是量子力学理论的对称性。

% [原文 p.68]
在具有连续规范群和连续 0-形式对称性的 4d 规范理论语境下，ABJ 反常由三角图刻画：其中两条顶点上是规范流，第三条顶点上是整体流。
\item \para{规范反常：} 当动力学规范场的规范变换不使其对路径积分的贡献保持不变时，就会出现这类反常。在这种情况下，所得的规范理论没有良好定义。在具有连续规范群的 4d 规范理论语境下，这类反常对应的三角图在三个顶点上都是规范流。
\item \para{'t Hooft 反常：} 在具有连续规范群的 4d 规范理论语境下，这类反常对应的三角图在三个顶点上都是整体流。
\end{itemize}
\end{notebox}

\begin{examplebox}{例 4.5：Maxwell 理论}
考虑 $d$ 维纯 $\U(1)$ 规范理论，同时打开电与磁高阶形式对称性的非平凡背景场，这相当于把 (4.14) 与 (4.15) 两项加入作用量。对电背景场作规范变换
\begin{equation}
  B^e_2\to B^e_2+\dd\Lambda^e_1\,, \tag{4.38}
\end{equation}
会使作用量改变
\begin{equation}
  S\to S+2\pi\int\Lambda^e_1\wedge\dd\star F\,. \tag{4.39}
\end{equation}
在磁背景存在时，$\star F$ 的运动方程被修改为
\begin{equation}
  \dd\star F=-\dd B^m_{d-2}\,. \tag{4.40}
\end{equation}
把它代入 (4.39)，我们看到配分函数按如下方式改变
\begin{equation}
  Z\left[B^e_2+\dd\Lambda^e_1,B^m_{d-2}\right]=\exp\left(-2\pi i\int\Lambda^e_1\wedge\dd B^m_{d-2}\right)\times Z\left[B^e_2,B^m_{d-2}\right]\,. \tag{4.41}
\end{equation}
由此可见，Maxwell 理论的电与磁高阶形式对称性之间存在混合 't Hooft 反常。

如果在上面的计算中交换电与磁的角色，我们也能看到这个反常。对磁背景场作规范变换
\begin{equation}
  B^m_{d-2}\to B^m_{d-2}+\dd\Lambda^m_{d-3}\,, \tag{4.42}
\end{equation}
在电背景存在时，会导致配分函数的如下反常变化
\begin{equation}
  Z\left[B^e_2,B^m_{d-2}+\dd\Lambda^m_{d-3}\right]=\exp\left(2\pi i\int\dd B^e_2\wedge\Lambda^m_{d-3}\right)\times Z\left[B^e_2,B^m_{d-2}\right]\,. \tag{4.43}
\end{equation}
\end{examplebox}

% [Tutor Note: p.68 式 (4.41) 可逐步核对：由 (4.39)，$\Delta S=2\pi\int\Lambda^e_1\wedge\dd\star F$；把 (4.40) 代入得 $\Delta S=-2\pi\int\Lambda^e_1\wedge\dd B^m_{d-2}$，故 $Z\to\ee^{\ii\Delta S}$ 恰给出 (4.41) 中的 $-2\pi i$ 相位。同法可验证 (4.43) 的符号为正：那里是 $\dd B^e_2$ 与 $\Lambda^m_{d-3}$ 配对，次数分别为 $3$ 与 $d-3$。]

% [原文 p.69]
\begin{figure}[htbp]
\centering
\includegraphics[width=0.62\textwidth]{fig13.png}
\figtitle{13}{'t Hooft 反常的两个例子。在上面的情形中，$U_{g'}$ 在 $U_g$ 下带电，因为把 $U_g$ 移过 $U_{g'}$ 会改变关联函数。在下面的情形中，可以类似地说，$U_{gg'}$ 与 $U_{g''}$ 之间的联结在 $U_{g'}$ 下带电。}
\end{figure}

\para{用拓扑算符描述。} 现在考虑离散对称性。上面我们讨论了背景场的规范变换等价于生成这些对称性的拓扑缺陷的形变。因此，当这类形变不是拓扑的、并且改变关联函数时，就出现 't Hooft 反常。这只有在拓扑缺陷或它们的联结彼此相交时才可能发生。见图 13。等价地，我们有下述陈述

\begin{statementbox}{陈述 4.1：'t Hooft 反常即带荷的拓扑缺陷}
一组离散对称性的 't Hooft 反常出现于：这些对称性的拓扑缺陷，或此类拓扑缺陷之间的联结，在该对称性下带有非平凡的荷。见图 13。
\end{statementbox}

添加抵消项对应于重新定义拓扑算符及其联结。

% [Tutor Note: p.69 图 13 上图“$U_{g'}$ 在 $U_g$ 下带电”的说法可用 $d=2$ 的 $\Z_2$ 0-形式对称性核对：把 $U_g$ 移过 $U_{g'}$ 得到的相位由两条线（或线与点）的交叉数给出，正是 $H^3(G,\U(1))$ 型反常的几何图像，与 p.74 的 (4.60) 相对应。]

% [原文 p.70]
\begin{examplebox}{例 4.6：离散规范理论}
考虑如 (4.18) 那样同时打开电与磁对称性背景场的离散（高阶形式）规范理论。执行例 4.5 中计算的离散版本，我们发现在这两种对称性之间存在混合 't Hooft 反常。作规范变换
\begin{equation}
  B^e_{p+1}\to B^e_{p+1}+\delta\Lambda^e_p\,, \tag{4.44}
\end{equation}
在存在 $B^m_{d-p}$ 时，配分函数按如下方式改变
\begin{equation}
  Z\left[B^e_{p+1}+\delta\Lambda^e_p,B^m_{d-p}\right]=\exp\left(2\pi i\,s(p)\int\Lambda^e_p\underset{\eta}{\cup}B^m_{d-p}\right)\times Z\left[B^e_{p+1},B^m_{d-p}\right]\,, \tag{4.45}
\end{equation}
其中我们定义了
\begin{equation}
  s(n):=(-1)^n\,,\qquad n\in\Z\,. \tag{4.46}
\end{equation}
类似地，作规范变换
\begin{equation}
  B^m_{d-p}\to B^m_{d-p}+\delta\Lambda^m_{d-p-1}\,, \tag{4.47}
\end{equation}
在存在 $B^e_p$ 时，配分函数按如下方式改变
\begin{equation}
  Z\left[B^e_{p+1},B^m_{d-p}+\delta\Lambda^m_{d-p-1}\right]=\exp\left(-2\pi i\int B^e_{p+1}\underset{\eta}{\cup}\Lambda^m_{d-p-1}\right)\times Z\left[B^e_{p+1},B^m_{d-p}\right]\,. \tag{4.48}
\end{equation}
这个 't Hooft 反常与下述事实有关：构成 $G^{(p)}$ 电 $p$-形式对称性的拓扑算符在 $\widehat{G}^{(p)}$ 磁 $(d-p-1)$-形式对称性下带电，反之亦然。规范变换 (4.47) 对应于移动磁对称性的拓扑算符，使它们扫出一条与上链 $\Lambda^m_{d-p-1}$ Poincaré 对偶的 $(p+1)$-链。由于背景 $B^e_{p+1}$ 的存在，时空内部插入了电对称性的拓扑算符网络。于是，在上述形变下，磁算符会穿过电算符，产生一个相位，它由电算符所占据的链与磁算符沿之形变的链之间的相交来描述。由 Poincaré 对偶，这些相位的乘积计算为
\begin{equation}
  \exp\left(-2\pi i\int B^e_{p+1}\underset{\eta}{\cup}\Lambda^m_{d-p-1}\right)\,. \tag{4.49}
\end{equation}
它正是上面讨论的配分函数的反常变化。
\end{examplebox}

% [Tutor Note: p.70 式 (4.48) 之前的正文写 “in the presence of $B^e_p$”，按 (4.44)(4.48) 的自洽性（此时电背景是 $B^e_{p+1}$）应为 $B^e_{p+1}$；已用 800 dpi 裁图确认原文印作 $B^e_p$，译文照录。]

% [Tutor Note: p.70 式 (4.49) 与 (4.48) 指数上的相位逐字相同，可作抄写核对；(4.45) 中引入的 $s(p)=(-1)^p$（见 (4.46)）是为离散情形中 $\cup$ 积的分次交换符号服务的约定：$p$ 的奇偶决定 (4.45) 的相位相对 (4.48) 的符号。]

\para{2d 中离散 0-形式对称性的 't Hooft 反常。} 利用上述用拓扑缺陷的移动来描述 't Hooft 反常的方式，我们可以完全刻画 $d=2$ 中 0-形式对称性可能的 't Hooft 反常。

% [原文 p.71]
\begin{figure}[htbp]
\centering
\includegraphics[width=0.35\textwidth]{fig14.png}
\figtitle{14}{与 2d 中 0-形式对称性相关联的拓扑局域算符与线算符的一种构型。}
\end{figure}

\begin{figure}[htbp]
\centering
\includegraphics[width=0.8\textwidth]{fig15.png}
\figtitle{15}{执行一次 F-移动会把关联函数改变一个 3-群上链 $\omega$ 倍。}
\end{figure}

考虑由离散群 $G^{(0)}$ 描述的 0-形式对称性。也就是说，我们有拓扑线算符
\begin{equation}
  L_g\,,\qquad g\in G^{(0)}\,. \tag{4.50}
\end{equation}
此外，在 $L_{g_1}$ 与 $L_{g_2}$ 汇合并转变为线算符 $L_{g_1g_2}$ 的联结处，选择拓扑局域算符
\begin{equation}
  \cO_{g_1,g_2}\,, \tag{4.51}
\end{equation}
见图 14。

考虑图 15 所示的拓扑线的移动，它有时被称为\textbf{缔合子}（associator）或\textbf{F-移动}（F-move）。这样的移动可能以如下复因子改变关联函数

% [Tutor Note: p.71 式 (4.50)(4.51) 建立了 2d 中 0-形式对称性的“线—结点”数据；下一片 (4.52) 的 $\omega(g_1,g_2,g_3)$ 正是三个线算符缔合时的相位，与 (4.55) 的五边形恒等式配套，可对照图 15 的 F-移动理解。]
